\reportsection{Motivation}
\label{sec:motivation}

\paragraph{The commitment layer is load-bearing, and sizing it is not the
user's job.}
Merkle commitments are the hash-based vector commitments practitioners actually
use: they rest only on a collision-resistant hash, which keeps them post-quantum
secure, and every interactive oracle proof in \textsf{Arkworks} is compiled
through one, so the commitment's guarantees become the compiled argument's
guarantees.  Configuring one correctly is not simple.  \textsf{ark-vc} fixes the
interface between the security definitions and interchangeable implementations,
\textsf{ark-mt} is the principal hash-based one, and a correct configuration
jointly sizes a digest width, a salt cardinality, and a domain-separation and
query-budget discipline against the target, tightly enough not to spend
efficiency it need not.  These widths move with field capacity and cannot be
hand-derived safely across fields, and an undersized choice still compiles and
passes every functional test.  Worse, a configuration cannot satisfy a property
the developer was never told it had to satisfy: knowing which claim a
configuration must support is itself a prerequisite to security, and the
infrastructure should hand the user that knowledge, because otherwise it is
hidden away and easily missed, with detrimental security consequences.

The definitions carry direct protocol meaning.  \emph{Binding} forces two
accepted views of one root to agree wherever they overlap, so a soundness
argument may treat committed answers as fixed; \emph{hiding} keeps the commitment
and its openings from revealing unopened entries, which is what a zero-knowledge
argument relies on.  This report scopes its claims to these two, with
extractability and equivocation named only as next steps.

\paragraph{Both defaults fail for small fields.}
The failure lives at the \code{MerkleHasher} trait, where the digest width and
the per-leaf salt are associated types that can be silently undersized with no
compile-time warning.  For BabyBear at $\secpar=128$, a one-element salt carries
only ${\sim}30$ bits, far below the entropy floor, so an adversary can recover a
hidden leaf and hiding is void.  A one-element digest publishes a root in a
${\sim}30$-bit space; an honest tree over $2^{20}$ leaves already makes roughly
$2$~million internal hash calls, far past the ${\sim}46{,}000$ birthday
threshold, and a deliberate collision confirms a binding break in $44$~seconds.
Neither failure implicates the hash function; both follow from
field-size-agnostic parameter choices.

\paragraph{Protocol use makes both failures worse, not safer.}
These numbers describe the commitment in isolation, but it is never deployed in
isolation.  In the BCS compilation \textsf{ark-vc} targets, which turns an
interactive oracle proof into a non-interactive argument by committing each
oracle message and answering the verifier's queries with authenticated openings,
the root is fixed into the argument transcript from the first message, and
binding enters the compiled soundness bound \emph{additively}, as
$\varepsilon_{\mathrm{PCP}}+\collisionBound(\digestbits,\querybudget)+\varepsilon$,
where $\varepsilon_{\mathrm{PCP}}$ is the soundness error of the underlying proof
system and $\varepsilon$ the residual slack the compilation adds.
Fiat--Shamir only tightens the digest requirement, since grinding lets the prover
hash many challenge nonces in search of a favourable transcript, and each attempt
is another oracle query, so the effective query budget inflates and the target
rises from roughly $2\secpar$ to $3\secpar$ bits.  Hiding has no protocol-level
rescue at all: the root sits in the transcript and is read before any later
message, so salt entropy cannot be traded for code distance, query count, or
round structure.  The fix is therefore not a better hash but a discipline that
sizes the parameters from the target and refuses the undersized ones.  The report
then follows it by property: binding (\cref{sec:binding}) and hiding
(\cref{sec:hiding}) carry their own definition, failure, and parameter,
\cref{sec:spine} unifies the two under one rule, and \cref{sec:cost} measures the
cost.
